\section{证据准入契约}
\label{sec:contract}
\subsection{状态与观察}
设工作区状态为
\begin{equation}
 S=(D,C,P,G,E,W),
\end{equation}
其中，$D$ 为设计声明，$C$ 为代码，$P$ 为适用策略，$G$ 为已观察的软件图，$E$ 为证据历史，$W$ 为工作项状态。其修订表示为 $\rho(S)=(\rho_C,\rho_D)$，其中设计分量可选。策略单独绑定，不是第三个修订字段。将设计纳入修订很有必要，因为即使代码字节未变，验收标准的变化也可能使先前结果不再适用。

为便于说明，将证据记录表示为
\begin{equation}
 e=(u,k,p,\rho,\sigma,a,o,\pi,\tau),
\end{equation}
其中，$u$ 为对象，$k$ 为证据类型，$p$ 为生产者，$\rho$ 为采集时的修订，$\sigma$ 为覆盖目标，$a$ 为权威性，$o$ 为结果，$\pi$ 为相关策略绑定，$\tau$ 为顺序标记或执行回执。该元组概括了分散在实现各处的检查，并不意味着所有子系统都采用同一种记录模式。

来源与权威性的作用不同。生产者标识记录某项观察由谁提供；权威性决定该入口有权确立哪些主张。自报的生产者名称或置信分数不能将外部证据转变为原生验证结果。同样，自动策略授权与明确的人工批准是不同的证据类型。如果把其中任何一组区别都替换成泛化的“成功”，就会改变决策的含义。

\subsection{适用性与生产者冲突}
对于请求 $q$，定义适用证据子集
\begin{align}
 E_q(S)&=\{e\in E:\App(e,S,q)\},\\
 \App(e,S,q)&=\operatorname{Complete}(e.\rho) \nonumber\\
 &\quad\land(e.\rho=\rho(S)) \nonumber\\
 &\quad\land\operatorname{Relevant}(e.\sigma,q) \nonumber\\
 &\quad\land\operatorname{Authorized}(e.a,P,q) \nonumber\\
 &\quad\land\operatorname{Compatible}(e.\pi,P,q).
 \label{eq:applicable}
\end{align}
各谓词的具体含义取决于操作类型。相关性用于纳入涉及请求目标的观察；当某条记录为某项具体要求提供支持时，再检查其确切覆盖范围。兼容性关注与该证据类型相关的策略条件。不完整的修订指纹不能确立与确切修订绑定的原生证据。为便于检查而保留这类记录，并不会使其具备适用性。

仅有适用性仍不足以解决记录之间的冲突。我们依据保留对象、证据类型、生产者和适用范围的键，对各组记录进行归约。只有满足相应的带类型回执规则与覆盖规则时，较新的记录才能取代较旧的记录。令 $L_q$ 表示归约后保留的记录。全局采用“以最新结果为准”的归约并不合适：生产者 $p_2$ 的通过结果不能抹去生产者 $p_1$ 的适用失败记录。

令 $Q$ 为非空的必需检查集合，$A$ 为额外配置的验收门控条件，$\bar S=(D,C,P,G,W)$ 为不含证据历史的状态。一条简化的准入规则为
\begin{align}
 \Accept(S,q) \iff{}& Q\ne\varnothing \nonumber\\
 &\land\ \forall t\in Q:\Pass(t,L_q) \nonumber\\
 &\land\ \nexists e\in L_q:\Fail_q(e) \nonumber\\
 &\land\ \operatorname{GatesSatisfied}(A,\bar S,L_q).
 \label{eq:accept}
\end{align}
这里，$\Pass(t,L_q)$ 要求适用的支持证据，其覆盖范围、类型、权威性与回执均需满足检查 $t$ 的要求；尚未补做的跳过检查不能提供这种支持。$\Fail_q(e)$ 表示请求策略下具有阻断作用的失败，并非泛指所有负面或提示性观察。已配置的门控条件仅通过 $L_q$ 访问证据。实现中区分了检查结果聚合与更大范围的协调核对和验收。式~\eqref{eq:accept} 将它们的逻辑依赖显式化，而不是用单一执行状态标志替代这些区别。

\begin{proposition}[加入无关记录时的保持性]
\label{prop:preservation}
固定 $\bar S$、请求 $q$、必需检查与门控条件，以及归约规则。加入不适用的证据不会改变准入结果。如果生产者 $p_1$ 的一条适用且具有阻断作用的失败记录仍被保留，则加入另一生产者 $p_2$ 的通过记录不能使准入条件变为真。
\end{proposition}
\begin{proof}
不适用记录在归约前即被排除，因此 $L_q$ 不变。检查与门控条件的其他参数均已固定，所以准入结果不变。不同生产者的记录具有不同的归约键，因此新增的通过记录不能取代该阻断失败记录。该失败记录使式~\eqref{eq:accept} 的第三个合取项为假。
\end{proof}

该命题只涉及指定的准入规则，并以输入真实、完整且绑定正确为前提。它不证明测试充分、生产者可信或实现已经得到形式化验证。尤其是，仅凭采集到的修订标签，不能证明某个进程在整个执行期间读取的都是不可变快照。这些假设区分了契约能够保证的性质与证据采集环节必须确立的事实。

\subsection{贯穿示例：一次变更}
表~\ref{tab:trace} 给出会话刷新示例的一条示意轨迹。标识符仅用于说明，不构成额外的实证任务。该轨迹揭示了两个彼此独立的有效性边界：编辑所使用的文件观察，以及验收所使用的修订绑定支持证据。

\begin{table}[t]
\small
\caption{一次变更中的证据状态转移示例。}
\label{tab:trace}
\begin{tabularx}{\columnwidth}{@{}p{0.19\columnwidth}X@{}}
\toprule
事件 & 对后续操作的影响\\
\midrule
读取 $h_0$ & 响应标识函数、原始函数体以及当前文件哈希。\\
其他工作者修改 & 文件状态变为 $h_1$；仍以 $h_0$ 为预期值的编辑必须刷新输入。\\
先前通过 & 当前修订为 $\rho_1$ 时，绑定到 $\rho_0$ 的证据仍作为历史记录保留。\\
$p_1$ 失败 & 当前失败记录在其对象与范围内保持可见。\\
$p_2$ 通过 & 来自另一生产者的通过记录不能取代 $p_1$ 的失败记录。\\
任务结束 & 完成操作更新工作项状态；验收仍需要适用的检查结果和门控条件。\\
\bottomrule
\end{tabularx}
\end{table}

源码摘录、测试通过记录和人工批准之间不存在一种通用排序，使其可以被视为可相互替代的置信度增量。它们支持的是不同谓词。因此，运行时同时公开观察结果和尚待满足的要求，使调用方能够有针对性地重新读取、重新验证或进行审查，而不必从一个未加区分的分数中推断就绪状态。
